% ================================================================
% Section: Challenges and Future Directions
% From main (2).tex — Section VIII + sections/conclusion.tex future work
% ================================================================

\section{Challenges and Future Directions}
\label{sec:challenges}

Despite remarkable progress, \glspl{slm} face persistent challenges that define the frontier of the field, and the agentic framing of Section~\ref{sec:agentic} adds a set of its own. We organize these into research clusters, closing with the empirical-verification agenda for this paper's consolidated evidence base.

\subsection{Hallucination Mitigation in Small Models}
\label{subsec:challenges_hallucination}

Hallucination, the confident generation of factually incorrect content, is an intrinsic failure mode of autoregressive language models, and it is more pronounced in smaller models \citep{zhang2023siren}. The core issue is information capacity: a 3B model's parameters can encode far fewer facts than a 175B model, and when queried about knowledge at the boundary of its training distribution, small models are more likely to confabulate plausible-sounding but incorrect responses.

Mitigation strategies include: (1)~\emph{\gls{rag}} \citep{lewis2020rag}, which anchors \gls{slm} generation to retrieved documents rather than parametric memory; (2)~\emph{uncertainty quantification} through sampling and ensemble methods; (3)~\emph{RLHF with truthfulness rewards} to penalize confident incorrect responses; and (4)~\emph{calibration post-processing} to align model confidence to actual accuracy.

\subsection{Reasoning Limitations Under Parameter Constraints}
\label{subsec:challenges_reasoning}

Complex multi-step reasoning, particularly in mathematics, formal logic, and scientific inference, improves roughly as a power law with parameter count \citep{kaplan2020scaling}. At the 1B--7B scale, \glspl{slm} struggle with problems that require holding many intermediate states simultaneously or with deep deductive chains.

Emerging directions include: (1)~\emph{Test-time scaling}: allocating additional inference compute (more generation tokens, repeated sampling, tree-of-thought search) to substitute for model capacity; DeepSeek-R1 \citep{deepseek2025r1} achieves frontier-class reasoning in its 7B distilled variant through extended chain-of-thought generation. (2)~\emph{Neuro-symbolic integration}: augmenting \glspl{slm} with symbolic solvers (SAT solvers, theorem provers, calculators) to handle formal reasoning externally. (3)~\emph{Speculative decoding} \citep{leviathan2023speculative} using a small draft \gls{slm} to propose tokens verified by a larger model, enabling faster inference from larger models when latency permits.

\subsection{Open Problems for Agentic SLM Systems}
\label{subsec:challenges_routing}

Section~\ref{sec:agentic} argues that \glspl{slm} are the appropriate default substrate for the class of agent calls hypothesis~H names; several problems stand between that conditional position and routine practice. \textbf{Difficulty estimation without an oracle:} a router must decide, cheaply and before generation, whether a given call exceeds the local model's competence; current approaches use learned win-prediction \citep{ong2024routellm} or sub-task decomposition \citep{division_of_thoughts_2025}, but calibration under distribution shift is unsolved. \textbf{Multi-objective routing:} the decision must jointly weigh accuracy, latency, energy, monetary cost, privacy sensitivity, and network availability, and the reported multilingual and safety gaps (Section~\ref{subsec:reported_safety}) mean a size-only router silently degrades non-English and safety-critical calls. \textbf{Handoff without a seam:} escalation from \gls{slm} to \gls{llm} mid-task must not produce a visible discontinuity in style or quality. \textbf{Agent-appropriate evaluation:} generalist leaderboards do not measure the tool-selection, argument-filling, and format-adherence skills agents actually use \citep{belcak2025slm_agentic}; function-calling leaderboards \citep{patil2025bfcl,zhang2024xlam} and interaction benchmarks \citep{yao2024taubench} are a start, but they report sub-10B models sparsely, and none reports cost per completed task (Section~\ref{subsubsec:agentic_ev_routing}). \textbf{Fleet management:} an agent built from many fine-tuned sub-3B specialists (Section~\ref{subsec:agentic_conversion}) raises questions of versioning, drift detection, and retraining cadence that single-model deployments do not face. Apple's privacy-preserving cloud escalation \citep{apple2024intelligence,gunter2025afm} shows the pattern is production-viable; generalizing it beyond a single vendor's stack is open.

\subsection{Hardware Co-Design and NPU Optimization}
\label{subsec:challenges_hardware}

The efficiency of \gls{slm} inference depends as much on hardware as on model architecture. Modern consumer hardware increasingly includes dedicated AI accelerators: Apple ANE, Qualcomm Hexagon \gls{npu} \citep{qualcomm2024ondevice}, Samsung Exynos \gls{npu}, and Intel \gls{npu}. These accelerators are optimized for specific operation types (INT4/INT8 matrix multiplication) and data layouts.

A critical open problem is \emph{hardware-model co-design}: developing \gls{slm} architectures whose attention patterns, FFN widths, and layer counts align with the data bus widths and tile sizes of target \glspl{npu}. \gls{gqa} head counts of 8 or 16 align well with SIMD widths on ARM Cortex-A processors, while head counts of 32 may be suboptimal. Future \glspl{slm} trained with hardware constraints as first-class objectives, rather than retrofitting quantization after the fact, will likely yield significant efficiency improvements.

\subsection{Ethical Considerations: Bias Amplification in Compressed Models}
\label{subsec:challenges_ethics}

Model compression introduces an underexplored ethical risk: \emph{bias amplification}. Pruning and quantization may disproportionately affect model performance on underrepresented demographic groups, dialects, and languages, as the model's remaining capacity may preferentially serve the majority distributions present in training data \citep{blodgett2020language}. \citet{kurz2025pruning_vs_quant} observe that quantization's advantage over pruning is most pronounced for multilingual tasks, suggesting that pruning disproportionately degrades non-English performance.

Research priorities include: (1)~\emph{fairness-aware compression} that includes demographic parity metrics as part of the compression objective; (2)~\emph{representation auditing} of compressed models across language, dialect, gender, and race; and (3)~\emph{minimum performance floors} for safety-critical applications, ensuring compressed models meet baseline accuracy thresholds for all user populations.

\subsection{Sustainable AI: SLMs as a Green AI Alternative}
\label{subsec:challenges_sustainability}

The energy consumption of training and serving frontier \glspl{llm} has become a significant environmental and economic concern \citep{strubell2019energy,patterson2022carbon}. Training GPT-3 emitted an estimated 552 metric tons of CO$_2$ equivalent \citep{patterson2021carbon}; training a model two orders of magnitude smaller from scratch consumes correspondingly less, and per-query inference energy is lower in proportion to the compute each query performs.

\subsubsection{What the Environmental Claim Covers}
\label{subsubsec:challenges_env_scope}

Claims of this kind are routinely made at a level of generality the accounting
does not support, and we withdraw one we made in an earlier draft. The statement
that 10,000 concurrent small-model instances on consumer hardware consume less
total energy than a single multi-GPU server is not one we can source or derive:
its truth depends entirely on utilization on both sides, and a server at high
batch utilization is considerably more energy-efficient per token than an idle
consumer device performing occasional single-stream inference. We remove the
claim rather than qualify it.

What can be said requires the boundary of the accounting to be stated. A complete
comparison of a small-model deployment against a large-model one has to include:
the training energy of the small model itself; the training energy of any
\emph{teacher} used to distill it, which for a distilled model can dominate and
which vendors do not report; the energy of the distillation, quantization, or
pruning pipeline that produced the deployed artifact; the embodied carbon of
manufacturing the devices inference runs on, amortized over their service life
and over the share of that life spent on inference; the marginal grid intensity
where each side runs, which differs by more than a factor of ten between regions;
and the energy of model updates over the deployment's lifetime, which for
on-device models means redistributing weights to every device rather than
updating one server.

Almost none of these terms is reported in the literature. What the field has
instead is per-query inference energy on the two sides, which is a single term in
a much longer sum, and which does favor the small model by one to two orders of
magnitude for a query both can answer. The defensible claim is therefore narrow:
\emph{for a task a small model completes correctly, serving it with the small
model uses substantially less inference energy}. Whether a wholesale shift to
small models reduces total system energy is a question no published accounting
answers, and the rebound effect, that cheaper inference induces more of it, cuts
against it. As AI regulation incorporates carbon disclosure requirements, this
accounting will have to be done properly, and the inputs for it do not currently
exist.

\subsection{Empirical-Verification Opportunities}
\label{subsec:challenges_verification}

Because this review synthesizes reported performance rather than producing original benchmark runs, the most immediately valuable follow-up work is targeted empirical verification of the reported claims that do the most work in practitioner decisions. Four empirical studies would disproportionately reduce the residual uncertainty in the evidence base reviewed here:

\begin{description}
    \item[Harness-consistent reruns of the headline \glspl{slm} on the headline benchmarks] A controlled re-run of the seven most-cited \glspl{slm} (Phi-4-mini, Phi-4, Qwen2.5-3B, Qwen2.5-7B, Gemma-3-4B, Llama 3.2-3B, Mistral-7B) on \gls{mmlu}, \gls{gsm8k}, HumanEval, and MBPP under a single \texttt{lm-evaluation-harness} configuration would anchor every vendor-reported figure to an independent reproduction.
    \item[Quantization-degradation measurement] A controlled comparison of FP16 vs.\ INT8 vs.\ INT4 accuracy across the seven headline \glspl{slm} and four benchmarks, using the \gls{gguf}, GPTQ, and \gls{awq} toolchains referenced in the vendor reports, would sharpen the 1--4-point degradation figure reported in the literature.
    \item[Edge-throughput reproduction] Timed runs of the same seven \glspl{slm} in their best-supported INT4 format on a Raspberry~Pi~5, Jetson Orin Nano, Snapdragon reference device, and recent laptop CPU would calibrate the reported throughput ranges against modern silicon.
    \item[QLoRA-gain reproduction] A controlled \gls{qlora} (rank 64, $\alpha = 128$) adaptation of the seven \glspl{slm} on SST-2, MNLI, ARC-C, \gls{gsm8k}, and SQuAD~2.0 would quantify the reported fine-tuning gains under a single adaptation recipe.

    \item[Cost per completed task, small-model against large-model agent] The measurement this survey most needs and could not find (Section~\ref{subsubsec:agentic_ev_routing}). Fix one agentic workload with a checkable success criterion, for instance a $\tau$-bench domain \citep{yao2024taubench} or a WebArena subset \citep{zhou2024webarena}. Run three configurations end to end: a frontier model throughout; a specialized \gls{slm} throughout; and a routed heterogeneous system. Report, per configuration, end-to-end task success, the number of model calls per completed task including retries, escalation rate, wall-clock latency, and total monetary and energy cost per \emph{completed} task rather than per call. This single study would move the economic premise of hypothesis~H from grade~\evC{} to grade~\evA{}, and it would settle whether per-token savings survive the retry and escalation overhead that per-call accounting hides.

    \item[Agent-specific evaluation of the sub-10B cohort] Run the seven headline \glspl{slm} on the Berkeley Function-Calling Leaderboard \citep{patil2025bfcl} and on $\tau$-bench, reporting the pass\textsuperscript{$k$} consistency metric alongside single-attempt accuracy, with and without constrained decoding \citep{willard2023guided}. This would fill the sparse-coverage gap identified in Table~\ref{tab:agentic_evidence} and would separate well-formedness failures from content failures, which current reporting conflates.
\end{description}

The first four studies are practical on a single modern GPU on the time scale of weeks. The last two require an agent harness and a modest API budget but no unusual resources, and they are the two that bear on the paper's central hypothesis rather than on its survey material. Together they would convert several consolidated claims from ``reported by sources'' to ``empirically verified,'' and would supply the first grade-\evA{} evidence in an area that currently has none. They do not replace the review; they strengthen the claims that practitioners depend on most.

\subsection{Open Research Directions Beyond Empirical Verification}
\label{subsec:challenges_future}

Beyond direct empirical verification, several promising research directions emerge from this review:

\begin{description}
    \item[Multilingual \gls{slm} evaluation] Extending the consolidated performance evidence to non-English benchmarks would reveal whether the English-centric performance hierarchy reported in the literature holds across languages, particularly for models with multilingual training data (Qwen2.5, Gemma~3, BLOOM).

    \item[On-device continual learning] Developing methods for incremental adaptation of deployed \glspl{slm} to new data without full retraining would enable personalization and domain adaptation at the edge.

    \item[SLM ensembles vs.\ single \gls{llm}] Investigating whether ensembles of specialized \glspl{slm} (one for math, one for code, one for general knowledge) with a lightweight router can match single-LLM performance at lower aggregate cost.

    \item[Multimodal \glspl{slm}] Evaluating vision-language \glspl{slm} such as Phi-3-vision, Gemma~3 with vision, and Gemma~3n \citep{google2025gemma3n} on multimodal benchmarks would extend the comparison to the rapidly growing multimodal inference domain.

    \item[Hybrid architectures] The emergence of non-transformer \gls{slm} architectures (SambaY \citep{pal2025sambay}, Gated Delta Networks \citep{yang2025gated_delta}) suggests that the transformer's dominance may not extend to the edge, where linear-time inference is particularly valuable.

    \item[Federated fine-tuning] Developing privacy-preserving fine-tuning protocols that aggregate gradient updates across edge devices would enable collaborative model improvement without centralized data collection.
\end{description}
